\begingroup
\small\setstretch{1.12}\setlength{\tabcolsep}{3pt}
\setlength{\LTleft}{\fill}\setlength{\LTright}{\fill}
\begin{longtable}{@{}>{\raggedright\arraybackslash}p{0.16\linewidth}>{\raggedright\arraybackslash}p{0.28\linewidth}>{\raggedright\arraybackslash}p{0.48\linewidth}@{}}
\caption{深度學習分類器之既有設定紀錄與更新狀態}\label{tab:deep-reference-settings}\\
\toprule
分類器 & 項目 & 保存內容 \\
\midrule\endfirsthead
\multicolumn{3}{c}{表~\thetable（續）}\\
\toprule
分類器 & 項目 & 保存內容 \\
\midrule\endhead
\midrule\multicolumn{3}{r}{續下頁}\\\endfoot
\bottomrule\endlastfoot
\multirow[t]{8}{=}{TabM} & 展示情境 & Old 1999--2011、Hybrid；由已保存驗證選參紀錄取值 \\*
 & \nolinkurl{d_block} & 512 \\*
 & \nolinkurl{dropout} & 0.2 \\*
 & \nolinkurl{k} & 32 \\*
 & \nolinkurl{lr} & 0.001 \\*
 & \nolinkurl{n_blocks} & 3 \\*
 & \nolinkurl{use_num_embeddings} & False \\*
 & \nolinkurl{weight_decay} & 1e-05 \\
\addlinespace[4pt]
TabICL & 選參與後端 & 原 CSV 未保存選定參數及可確認的後端版本 \\
\addlinespace[4pt]
兩分類器 & 事件標籤 & 待以相同年份、驗證規則與種子重跑後更新 \\
\end{longtable}
\noindent{\footnotesize (資料來源：本研究既有原始結果欄位與執行入口核對)\par TabM 僅展示預先固定情境的一組已選參數，不代表所有邊界皆使用相同設定；不由所選值反推完整歷史搜尋預算。未保存的 TabICL 設定不依目錄名稱或測試分數補造。}
\endgroup
